\documentclass[10pt,a4paper]{amsart}
\usepackage[margin=19mm]{geometry}
\usepackage{amsmath,amssymb,mathtools}
\usepackage[scr=rsfs,cal=euler]{mathalfa}
\usepackage{xfrac}
\usepackage[unicode,hidelinks,bookmarks=false]{hyperref}
\newcommand{\ie}{i.e.\ }
\newcommand{\cf}{cf.\ }
\newcommand{\resp}{respectively\ }
\newcommand{\Subsection}{\subsection}
\providecommand{\nobreakdash}{\nobreak\hbox{-}}
\setlength{\emergencystretch}{2em}
\multlinegap0pt
\usepackage{amssymb}
\usepackage{tcolorbox}
\usepackage{tikz}
\usepackage{xcolor}
\newtheorem{theorem}{Theorem}
\DeclareMathOperator{\diomega}{\ora \omega}
\newcommand{\norm}[1]{\left|#1\right|}
\newcommand{\ora}[1]{\overrightarrow{#1}}
\title{A Note on the Complexity of Directed Clique}
\author{Grzegorz Gutowski, Miko{\l}aj Rams}
\date{Source: arXiv:2602.11773v1 (CC BY 4.0)}
\begin{document}
\maketitle

\noindent\emph{Source and licence.} A Note on the Complexity of Directed Clique. Version arXiv:2602.11773v1 (CC BY 4.0), distributed by arXiv under the Creative Commons Attribution 4.0 International licence, http://creativecommons.org/licenses/by/4.0/, as recorded in the arXiv licence metadata for this exact version and shown on the abstract page https://arxiv.org/abs/2602.11773v1. Full licence notice: \texttt{licenses/arxiv-2602.11773-CC-BY-4.0.txt}.\par\smallskip

\noindent\textit{Editorial scope.} Counted unit: the article's theorem thm:upper\_bound (source0.tex lines 641--643). The article's complete original proof of it is delivered with it (source0.tex lines 645--664, one \texttt{proof} environment of 20 lines). No mathematics is written, repaired or re-derived: the statement and the proof are the article's own lines, in the article's own notation, and no retained line is substituted or re-flowed.  No further source-local context is needed: every cross-reference of every retained line resolves inside the two delivered spans.  Nothing was omitted from any retained span, so the package carries no omission record.  No retained line contains a citation, so the wrapper carries no bibliography and no bibliography entry.  No retained line contains a figure environment, an image-inclusion command or drawing code, so no image is delivered and the package declares no asset dependency.  Editorial formatting only, in addition: an AMS \texttt{amsart} preamble that restates the article's own declarations and macros used by the retained lines, namely the environments \texttt{theorem} on separate counters, together with the standard packages those lines need.  The retained environments print in this document as Theorem 1 in order of appearance, whereas the article prints the same items with the numbering of its own full document.  The complete native source of the article as distributed in the arXiv e-print for this exact version is bundled unchanged as \texttt{sources/194488/source0.tex}.
\begin{theorem}\label{thm:upper_bound}
    For every tournament \(T\) with \(\diomega(T) \ge t \ge 1\), we have \(\norm{V(T)} \geq {{t+1}\choose{2}}\).
\end{theorem}

\begin{proof}
    The proof goes by the induction on $t$.
    For the base case $t=1$, we obviously have that the tournament needs to have at least $1={{1+1} \choose 2}$ vertex.
    Now, for the induction step, let $t\ge2$, and assume that the statement of the theorem holds for $t-1$.

    Let $T$ be a tournament with $\diomega(T) \ge t$.
    As $\diomega(T) \ge t$, then for every linear order of $V(T)$, the backedge graph contains a clique $K_t$.
    Thus, there is at least one subset $U \subseteq V(G)$ of vertices that induces a copy of a transitive tournament $TT_t$.
    Let $U' = V(G)\setminus U$ be the complement of $U$, $T' = T[U']$ be the tournament induced by $U'$, $t' = \diomega(T')$ be the directed clique number of $T'$, and $\ll'$ be some linear order of $U'$ with $\omega(T'^{\ll'})=t'$.

    As $U$ induces a transitive subtournament in $T$, we can define order $\ll''$ on $U$ such that for any $x,y \in U$ we have $x \ll'' y$ if and only if $(x,y) \in A(T)$.
    Consider a linear order $\ll$ on $V(G)$ that first puts all vertices in $U'$ in order $\ll'$ and then puts all the vertices in $U$ in order $\ll''$.
    As $\diomega(T) \ge t$, there is at least one copy of a clique $K_t$ in the backedge graph $T^\ll$.
    We have that there are no edges between vertices in $U$ in the backedge graph $T^\ll$.
    Thus, any clique in the backedge graph $T^\ll$ can include at most one vertex in $U$, and we have $t' \ge t-1$.
    By the induction hypothesis, we get $\norm{U'} \ge {t\choose2}$, and
    \[
        \norm{V(G)} = \norm{U'} + \norm{U} \ge {t\choose2} + t = {{t+1}\choose2}\text{.}
    \]
\end{proof}

\end{document}
